\section*{第 \ref{ch:g10:refraction} 章 --- 光的折射}

\begin{solution}{exo:g10:refraction:1}
角度自法线量起：入射角为
$90^\circ - 25^\circ = 65^\circ$，故反射角也为
$65^\circ$。入射光束与反射光束的
夹角为 $2 \times 65^\circ = 130^\circ$。
\end{solution}

\begin{solution}{exo:g10:refraction:2}
\emph{1.} 由 $v = c/n$，在水中
$v = \qty{3.00e8}{m/s}/1.33 = \qty{2.26e8}{m/s}$；在金刚石中
$v = \qty{3.00e8}{m/s}/2.42 = \qty{1.24e8}{m/s}$。

\emph{2.} $n = \dfrac{c}{v} = \dfrac{\qty{3.00e8}{m/s}}
{\qty{1.97e8}{m/s}} = 1.52$：玻璃。
\end{solution}

\begin{solution}{exo:g10:refraction:3}
$\sin i_2 = \dfrac{\sin 40^\circ}{1.50} = \dfrac{0.643}{1.50} = 0.429$，
故 $i_2 = 25.4^\circ$。由光路可逆性，玻璃内
以 $25.4^\circ$ 射到界面上的光线，
恰好以 $40^\circ$ 的角度射入空气。
\end{solution}

\begin{solution}{exo:g10:refraction:4}
$\sin i_2 = \dfrac{\sin 55^\circ}{1.33} = \dfrac{0.819}{1.33} = 0.616$，
故 $i_2 = 38.0^\circ$。光线\emph{向法线}偏折，原因是
水为光速较慢的介质（$n_2 > n_1$）。
\end{solution}

\begin{solution}{exo:g10:refraction:5}
$n = \dfrac{\sin 45^\circ}{\sin 32^\circ} = \dfrac{0.707}{0.530}
= 1.33$：该液体为水。
\end{solution}

\begin{solution}{exo:g10:refraction:6}
\emph{1.} 五个比值分别为
\[
\frac{0.342}{0.233} = 1.47,\quad
\frac{0.500}{0.334} = 1.50,\quad
\frac{0.643}{0.431} = 1.49,\quad
\frac{0.766}{0.515} = 1.49,\quad
\frac{0.866}{0.581} = 1.49 :
\]
在测量误差范围内为常量 --- 斯涅耳定律。

\emph{2.} $n \approx 1.49$：有机玻璃。

\emph{3.} 斯涅耳定律给出 $\sin i_2 = \sin i_1 / n$，即正弦对正弦，
数据须落在过原点的一条直线上，
直尺一看即可检验。若以角度对角度作图，数据落在
一条曲线上；小角度时，曲线看似与直线相近，
这恰好蒙蔽了托勒密。
\end{solution}

\begin{solution}{exo:g10:refraction:7}
$\sin i_c = n_2/n_1$，取 $n_2 = 1.00$：
水：$\sin i_c = 1/1.33 = 0.752$，$i_c = 48.8^\circ$；
玻璃：$\sin i_c = 1/1.50 = 0.667$，$i_c = 41.8^\circ$；
金刚石：$\sin i_c = 1/2.42 = 0.413$，$i_c = 24.4^\circ$。

由空气射入水中时，$\sin i_2 = \sin i_1/1.33 \leq 1/1.33 < 1$，对
任意入射角都成立：总存在折射角，因此不会出现任何特殊情形 ---
全反射需要 $n_1 > n_2$。
\end{solution}

\begin{solution}{exo:g10:refraction:8}
\emph{1.} $d' = \dfrac{d}{n} = \dfrac{\qty{2.0}{m}}{1.33}
= \qty{1.5}{m}$.

\emph{2.} 鱼的视深为 $\qty{40}{cm}/1.33 = \qty{30}{cm}$：
像位于鱼的\emph{上方}，故苍鹭必须击向\emph{下方}，
即比所见之像更低的位置。苍鹭不受斯涅耳定律拖累，
靠反复试错学会此事。
\end{solution}

\begin{solution}{exo:g10:refraction:9}
\emph{1.} $\sin r = \dfrac{\sin 50^\circ}{1.50} = \dfrac{0.766}{1.50}
= 0.511$，故 $r = 30.7^\circ$。

\emph{2.} 光线由玻璃射向空气，入射角为 $30.7^\circ$：
$\sin i_2 = 1.50 \times 0.511 = 0.766$，故 $i_2 = 50^\circ$ ---
出射角等于入射角。

\emph{3.} 两次折射相互抵消：光线离开时
与原方向\emph{平行}，仅发生侧向平移。
这就是窗玻璃不会使景物变形的原因，
只会（难以察觉地）使其平移。
\end{solution}

\begin{solution}{exo:g10:refraction:10}
$\sin 55^\circ = 0.819$。红光：
$\sin i_2 = 0.819/1.61 = 0.509$，$i_2 = 30.6^\circ$。蓝光：
$\sin i_2 = 0.819/1.66 = 0.493$，$i_2 = 29.6^\circ$。两种颜色
相距 $1.0^\circ$；蓝光折射率更大，偏折更多，
最终更靠近法线。
\end{solution}

\begin{solution}{exo:g10:refraction:11}
\emph{1.} $\sin i_c = \dfrac{1.46}{1.48} = 0.986$，故
$i_c = 80.6^\circ$。

\emph{2.} 与法线成 $i_c$ 角的光线，与轴线所成角为 $90^\circ - i_c$，
故沿轴向长 $L$ 的路径长为 $\dfrac{L}{\sin i_c}$：
\[
\qty{1000}{m} \times \frac{1.48}{1.46} = \qty{1013.7}{m},
\]
比轴向光线约多 \qty{14}{m}。
\end{solution}

\begin{solution}{exo:g10:refraction:12}
\emph{1.} 从地平线方向掠来的光，入射角 $i_1$ 接近
$90^\circ$，其折射满足 $\sin i_2 = \sin 90^\circ / 1.33 = 0.752$，
即 $i_2 = 48.8^\circ$ --- 恰为临界角，
这正是光路可逆性的要求。因此，来自整个天空的光线
都落在半顶角为 $48.8^\circ$ 的圆锥之内，方能到达眼睛。

\emph{2.} 圆盘的边缘，就是该圆锥与水面相交之处：
\[
r = d \tan i_c = \qty{2.0}{m} \times \tan 48.8^\circ
= \qty{2.0}{m} \times 1.14 = \qty{2.3}{m}.
\]

\emph{3.} 圆盘之外，来自水下的光在水面处的入射角超过
临界角，发生全反射：潜水者看到的是一面银白色的
镜子，映出池底。
\end{solution}

\begin{solution}{exo:g10:refraction:13}
\emph{1.} 红光：$\sin i_2 = 1.51 \times \sin 30^\circ = 0.755$，故
$i_2 = 49.0^\circ$。蓝光：$\sin i_2 = 1.53 \times 0.500 = 0.765$，故
$i_2 = 49.9^\circ$。白光光束的散开角为
$49.9^\circ - 49.0^\circ = 0.9^\circ$ --- 仅一次折射便形成
棱镜式光谱。

\emph{2.} 入射角为 $45^\circ$ 时：$\sin i_2 = 1.51 \times 0.707 = 1.07 > 1$
（蓝光为 $1.53 \times 0.707 = 1.08$）。两种颜色的光均不存在折射角，
其临界角为
$\arcsin(1/1.51) = 41.5^\circ < 45^\circ$，故长面
对整束光发生全反射。
\end{solution}

\begin{solution}{exo:g10:refraction:14}
\emph{1.} 玻璃与空气界面的临界角：$\sin i_c = 1/1.50$，故
$i_c = 41.8^\circ < 45^\circ$，光束发生全反射 --- $100\%$
的光被反射，无需镀膜，因而不会变暗。金属镜子
每次反射损失百分之几，潜望镜则需两次反射。

\emph{2.} 玻璃与水界面：$\sin i_c = \dfrac{1.33}{1.50} = 0.887$，故
$i_c = 62.5^\circ$。此时 $45^\circ < i_c$，入射角低于
临界角，大部分光折射进入水中，于是
镜子 --- 因而所成的像 --- 随之失效。

\emph{3.} $\sin i_2 = \dfrac{1.50}{1.33} \times \sin 45^\circ
= 1.128 \times 0.707 = 0.798$，故射入水中时 $i_2 = 52.9^\circ$。
\end{solution}

\begin{solution}{exo:g10:refraction:15}
\emph{1.} $v = \dfrac{\qty{3.00e8}{m/s}}{1.47} = \qty{2.04e8}{m/s}$，故
\[
t = \frac{\qty{6.2e6}{m}}{\qty{2.04e8}{m/s}} = \qty{30.4}{ms} \approx
\qty{30}{ms}.
\]

\emph{2.} 无线电波：$t = \qty{6.2e6}{m}/\qty{3.00e8}{m/s} =
\qty{20.7}{ms}$。玻璃造成的时间为 $30.4 - 20.7 = \qty{9.7}{ms}$，
单程如此，往返约 \qty{19}{ms}。

\emph{3.} 降低折射率 --- 空心光纤使光在空气中传播
（$n \approx 1.00$），微波链路经大气层传输亦然，
或缩短路径：将光缆铺设得更接近
两城之间的大圆航线。两种方案实际都有产品出售，
价格以每毫秒计。
\end{solution}

\begin{solution}{pb:g10:refraction:1}
\textbf{1.} $v_{\text{水}} = \qty{3.00e8}{m/s}/1.33 =
\qty{2.26e8}{m/s}$；$v_{\text{金刚石}} = \qty{3.00e8}{m/s}/2.42 =
\qty{1.24e8}{m/s}$ --- 金刚石使光速减慢的倍数为 $n = 2.42$。

\textbf{2.} $\sin i_2 = \sin 40^\circ / 1.33 = 0.643/1.33 = 0.483$，故
$i_2 = 28.9^\circ$。

\textbf{3.} $\sin i_2 = 1.33 \times \sin 28.9^\circ = 1.33 \times
0.483 = 0.643$，故 $i_2 = 40^\circ$：光沿原路返回 ---
即光线的可逆性。

\textbf{4.} $\sin i_c = 1/1.33 = 0.752$，故 $i_c = 48.8^\circ$。入射角为
$60^\circ > i_c$ 时，斯涅耳定律要求 $\sin i_2 = 1.33 \sin
60^\circ = 1.15 > 1$，不存在折射光线：水面发生全反射，
光线被反射回水中。

\textbf{5.} $v = \qty{3.00e8}{m/s}/1.0003 = \qty{2.999e8}{m/s}$：光
在空气中的速率比真空中慢 $0.03\%$，远低于我们
测角的精度，故取 $n_{\text{空气}} = 1.00$ 并无妨碍。

\textbf{6.} $d' = \qty{3.0}{m}/1.33 = \qty{2.3}{m}$：池水
看起来比实际深度浅 \qty{70}{cm}。

\textbf{7.} 从硬币射出的两条光线在水面发生折射，
\emph{远离}法线偏折，因此它们在水面上方发散得
比水下更陡。眼睛将两条出射光线沿直线向后延长，
它们相交于硬币\emph{上方}，这一交点
即为眼睛所确定的像的位置。

\textbf{8.} $r = d \tan i_c = \qty{3.0}{m} \times \tan 48.8^\circ =
\qty{3.0}{m} \times 1.14 = \qty{3.4}{m}$：直径约 \qty{6.9}{m}
的圆盘即可容纳整个天空。

\textbf{9.} 入射角为 $30^\circ$ 时：$\sin i_2 = 1.33 \times 0.500 = 0.665$，光线
以 $41.7^\circ$ 射出。入射角为 $45^\circ$ 时：$\sin i_2 = 1.33 \times
0.707 = 0.940$，光线以 $70.1^\circ$ 射出，几乎沿水面掠出。入射角为
$50^\circ$ 时：$1.33 \times 0.766 = 1.02 > 1$ --- 发生
全反射，光线重新射回池中。

\textbf{10.} 像位于真实鱼的上方
（见\cref{prop:g10:refraction:depth}），故应瞄准其\emph{下方}。
是的：鱼同样看到渔夫发生了偏移，在斯涅耳窗之内
被挤向竖直方向 --- 双方各自看到的对方，都不在对方的真实位置上。

\textbf{11.} 金刚石：$\sin i_c = 1/2.42 = 0.413$，于是
$i_c = 24.4^\circ$，玻璃则为 $41.1^\circ$：金刚石的逸出锥
锥宽还不到玻璃的一半。

\textbf{12.} 在金刚石中，$30^\circ > 24.4^\circ$，光线发生全
反射，光线留在宝石内。在玻璃中，
$30^\circ < 41.1^\circ$：光线从背面漏出。明亮式切割的刻面
倾斜角经过设计，使射入的光几乎总射到它们上，
入射角超过 $24.4^\circ$，光线不断反射，直至向上射出，朝向
眼睛——玻璃仿制品的逸出锥宽，光只会直接穿过。
线直接穿过。

\textbf{13.} $\sin 45^\circ = 0.707$。红光：$0.707/2.41 = 0.293$，于是
$i_2 = 17.1^\circ$。蓝光：$0.707/2.45 = 0.289$，于是 $i_2 = 16.8^\circ$。
入射时的色散约为 $0.3^\circ$。

\textbf{14.} $\sin 17.0^\circ = 0.292$。红光：$\sin i_2 = 2.41 \times
0.292 = 0.705$，于是 $i_2 = 44.8^\circ$。蓝光：$\sin i_2 = 2.45 \times
0.292 = 0.716$，于是 $i_2 = 45.8^\circ$。此时光扇的张角为 $1.0^\circ$，
出射折射在临界角附近进行，把入射处的
$0.3^\circ$ 色散拉大了约三倍。这种拉开的彩色光束
即金刚石的\emph{火彩}。

\textbf{15.} 金刚石--水：$\sin i_c = 1.33/2.42 = 0.550$，于是
$i_c = 33.3^\circ$。那条 $30^\circ$ 光线原本在刻面朝向空气时被困住，
现在可以逸出（$30^\circ < 33.3^\circ$），进入水膜：每一滴
水珠、指纹或油膜，都会提高外侧的折射率、扩大
逸出锥，并削弱闪光 --- 这正是珠宝商要用擦拭布的缘故。

\textbf{16.} $\sin i_c = \dfrac{1.46}{1.48} = 0.986$，于是
$i_c = 80.6^\circ$。

\textbf{17.} 全反射发生在\emph{表面处}，
因此表面必须保持完好：在裸露的光纤上，每一粒
灰尘、划痕、水滴或接触的手指，都会提高该处外侧的折射率，
并使光线泄漏出去。包层把全反射的界面
埋入洁净的玻璃之中，使其不受任何触碰 --- 另外，
还能从力学上保护细如发丝的纤芯。

\textbf{18.} $v = \qty{3.00e8}{m/s}/1.48 = \qty{2.03e8}{m/s}$，于是
\[
t = \frac{\qty{6.0e6}{m}}{\qty{2.03e8}{m/s}} = \qty{29.6}{ms} \approx
\qty{30}{ms}.
\]

\textbf{19.} 路径因子 $1/\sin i_c = 1.48/1.46 = 1.0137$：最坏情况下
$\qty{6000}{km} \times 0.0137 \approx \qty{82}{km}$ 的额外玻璃路程，
多耗时 $\qty{8.2e4}{m} / \qty{2.03e8}{m/s} \approx \qty{0.4}{ms}$ ---
之字形路径几乎不增加时间。

\textbf{20.} 往返延迟：$2 \times (29.6 + 0.4) \approx \qty{60}{ms}$。
吞吐量：每秒约可传送 $\dfrac{10^{13}}{2.4 \times 10^8} \approx 40\,000$
份此书的副本。一句话概括：借助全反射将光束困在玻璃内，
斯涅耳定律便能每秒把一座图书馆送过大洋，
来回仅需六十毫秒。
\end{solution}
